\documentclass[11pt,a4paper]{article}

\usepackage[utf8]{inputenc}
\usepackage[T1]{fontenc}
\usepackage{lmodern}
\usepackage[french]{babel}
\usepackage[a4paper,margin=2.15cm]{geometry}
\usepackage{longtable}
\usepackage{booktabs}
\usepackage{array}
\usepackage{xcolor}
\usepackage{hyperref}

\definecolor{zeloGreen}{HTML}{3B8F2F}
\definecolor{zeloSoft}{HTML}{F2F8EF}
\definecolor{zeloInk}{HTML}{1F2A1F}
\definecolor{riskRed}{HTML}{A83232}
\definecolor{riskOrange}{HTML}{B86B00}

\hypersetup{
  colorlinks=true,
  linkcolor=zeloGreen,
  urlcolor=zeloGreen
}

\setlength{\parskip}{0.55em}
\setlength{\parindent}{0pt}
\renewcommand{\arraystretch}{1.22}
\newcommand{\code}[1]{\texttt{\detokenize{#1}}}
\newcommand{\riskpzero}{\textcolor{riskRed}{\textbf{P0}}}
\newcommand{\riskpone}{\textcolor{riskOrange}{\textbf{P1}}}
\newcommand{\statusok}{\textcolor{zeloGreen}{\textbf{validé}}}

\title{\textbf{Zelo AI -- Correctif P0 d'isolation des sessions client}\\
\large Note technique de sécurité, validation multi-compte et impact conformité}
\author{Projet Zelo AI}
\date{18 juin 2026}

\begin{document}
\maketitle

\begin{center}
\fcolorbox{zeloGreen}{zeloSoft}{
  \begin{minipage}{0.92\linewidth}
  \textbf{Synthèse.} Une revue des logs et du code a mis en évidence un
  défaut critique d'isolation entre utilisateurs web : plusieurs sessions
  pouvaient partager le même identifiant technique \code{local-dev}. Le risque
  était qu'un panier ou une commande utilise le token Swissdata du mauvais
  compte. Le correctif remplace cette confiance dans le frontend par une
  session serveur opaque, transmise par cookie \code{HttpOnly}, et impose que
  les routes panier/chat/commande dérivent toujours l'identité depuis la
  session authentifiée côté serveur. Le correctif est ensuite étendu au canal
  WhatsApp via une authentification passwordless par code e-mail et une session
  serveur distincte par numéro.
  \end{minipage}
}
\end{center}

\tableofcontents
\newpage

\section{Objectif du document}

Ce document complète le livrable Phase 4 sécurité et conformité. Il décrit en
détail un correctif \riskpzero{} appliqué à Zelo AI après détection d'un risque
de confusion de sessions client.

L'objectif est double :
\begin{itemize}
  \item expliquer précisément pourquoi l'incident pouvait arriver;
  \item documenter la mitigation retenue selon les bonnes pratiques de
  cybersécurité : séparation des sessions, moindre confiance dans le client,
  réduction des données sensibles en logs et tests multi-utilisateurs.
\end{itemize}

Le document ne prétend pas être une certification externe. Il constitue une
preuve interne d'analyse, de correction et de validation, utile pour la Phase 4
du projet et pour les exigences nLPD/GDPR liées à la confidentialité et à
l'intégrité des données client.

\section{Constat initial}

\subsection{Symptôme observé}

Lors d'un test sur l'environnement déployé, une commande passée depuis un compte
client distinct a généré une confirmation envoyée à une autre adresse e-mail.
Des observations antérieures indiquaient aussi que des paniers pouvaient parfois
sembler partagés ou contaminés entre utilisateurs.

Ce type de symptôme est critique, car il touche directement :
\begin{itemize}
  \item la confidentialité du panier;
  \item l'intégrité de la commande;
  \item l'identité du compte facturé ou notifié;
  \item la confiance utilisateur dans l'assistant.
\end{itemize}

\subsection{Racine technique}

L'interface web utilisait historiquement une valeur fixe \code{local-dev} comme
identifiant utilisateur technique. Cette valeur était envoyée par le navigateur
dans les appels suivants :
\begin{itemize}
  \item login;
  \item chat;
  \item affichage panier;
  \item mutations panier;
  \item passage de commande.
\end{itemize}

Côté backend, les tokens Swissdata étaient stockés en mémoire par clé
\code{userRef}. Quand plusieurs utilisateurs web partageaient le même
\code{userRef}, le dernier login pouvait remplacer le token précédemment stocké.
Une commande ultérieure pouvait donc être exécutée avec le mauvais token
Swissdata.

\begin{longtable}{p{0.28\linewidth}p{0.66\linewidth}}
\toprule
\textbf{Élément} & \textbf{Problème avant correctif} \\
\midrule
Identité frontend & Le navigateur pouvait envoyer \code{userRef="local-dev"}. \\
Stockage token & Le token Swissdata était associé à cette clé partagée. \\
Isolation panier & Deux comptes web pouvaient se retrouver sur la même clé runtime. \\
Passage de commande & Le compte utilisé dépendait du token stocké en dernier sous la clé partagée. \\
Audit sécurité & Les logs confirmaient que des appels web utilisaient encore \code{local-dev}. \\
\bottomrule
\end{longtable}

\section{Classification du risque}

\subsection{Sévérité}

Le risque est classé \riskpzero{} parce qu'il peut provoquer une action
irréversible ou sensible sur le mauvais compte client :
\begin{itemize}
  \item commande passée avec le mauvais client Swissdata;
  \item confirmation envoyée au mauvais destinataire;
  \item panier d'un utilisateur visible ou modifiable depuis une autre session;
  \item violation potentielle de confidentialité et d'intégrité des données.
\end{itemize}

\subsection{Lien avec OWASP et bonnes pratiques}

Le problème correspond principalement à deux familles de risques :
\begin{itemize}
  \item \textbf{Broken Authentication / Session Management} : l'identité n'était
  pas suffisamment portée par une session serveur isolée.
  \item \textbf{Broken Object Level Authorization} : une action panier ou
  commande pouvait être appliquée au mauvais objet utilisateur si la clé runtime
  était partagée.
\end{itemize}

La mitigation suit une règle simple : le backend ne doit jamais faire confiance
à une identité utilisateur choisie par le client pour une action mutante.

\section{Design du correctif}

\subsection{Principe retenu}

Le correctif introduit une session web opaque créée côté serveur au moment du
login. Cette session devient la seule source d'identité pour les routes web
authentifiées.

\begin{quote}
Le navigateur peut demander une action, mais il ne choisit plus le compte
Swissdata sur lequel cette action s'exécute.
\end{quote}

\subsection{Flux corrigé}

\begin{verbatim}
POST /ai/login
  -> authentification Swissdata avec username/password
  -> création session serveur aléatoire
  -> stockage token sous userRef interne web:<sessionId>
  -> cookie HttpOnly envoyé au navigateur

POST /ai/chat
GET  /ai/basket
POST /ai/logout
GET  /ai/history-suggestions
  -> lecture cookie
  -> récupération session serveur
  -> récupération userRef interne
  -> action Swissdata avec le bon token
\end{verbatim}

\subsection{Avant / après}

\begin{longtable}{p{0.24\linewidth}p{0.34\linewidth}p{0.34\linewidth}}
\toprule
\textbf{Point} & \textbf{Avant} & \textbf{Après} \\
\midrule
Principal utilisateur & Fourni par le frontend via \code{userRef}. & Dérivé côté serveur depuis un cookie de session. \\
Valeur web typique & \code{local-dev}, partagée. & \code{web:<sessionId>}, aléatoire et par session. \\
Route chat/panier & Acceptait l'identité du client. & Ignore l'identité envoyée par le client. \\
Session expirée & Risque de comportement ambigu selon token mémoire. & Réponse \code{401} si aucune session valide. \\
Logout & Déconnexion frontend seulement ou partielle. & Suppression session serveur, token Swissdata et cookie. \\
Logs & Certains logs pouvaient exposer trop de détail panier. & Logs réduits et redaction de données sensibles. \\
\bottomrule
\end{longtable}

\section{Implémentation technique}

\subsection{Session serveur}

Le fichier \code{src/swissdata/session.ts} a été étendu pour gérer des sessions
web opaques :
\begin{itemize}
  \item génération d'un identifiant aléatoire via cryptographie Node;
  \item mapping \code{sessionId -> userRef};
  \item préfixe interne \code{web:} pour distinguer les sessions navigateur;
  \item stockage du username authentifié, de la date de création et de mise à
  jour;
  \item expiration et suppression des sessions;
  \item fonction de masquage pour les logs.
\end{itemize}

Cette approche garde l'état de session en mémoire, ce qui est suffisant pour le
prototype et le déploiement actuel mono-instance. Une section de limites détaille
les améliorations nécessaires en cas de scaling horizontal.

\subsection{Authentification et cookie}

Le fichier \code{src/routes/ai.routes.ts} introduit une couche d'authentification
web :
\begin{itemize}
  \item lecture contrôlée du cookie de session;
  \item helper \code{authenticatedUser} pour récupérer le principal courant;
  \item helper \code{requireAuthenticatedUser} pour bloquer les routes sensibles;
  \item cookie \code{HttpOnly};
  \item \code{SameSite=Lax};
  \item \code{Secure} en HTTPS ou en production;
  \item compatibilité locale sans \code{Secure} sur \code{localhost}.
\end{itemize}

Le cookie ne contient pas le token Swissdata. Il contient seulement un
identifiant opaque qui n'a de sens que pour le serveur Zelo AI.

\subsection{Routes protégées}

Les routes web sensibles ne prennent plus leur identité depuis le body ou la
query string :
\begin{itemize}
  \item \code{/ai/chat};
  \item \code{/ai/basket};
  \item \code{/ai/history-suggestions};
  \item \code{/ai/build-meal-search};
  \item \code{/ai/logout};
  \item \code{/ai/whoami}.
\end{itemize}

Même si un client envoie volontairement \code{userRef="local-dev"} dans le
payload, le backend l'ignore pour les actions authentifiées et utilise la
session serveur.

\subsection{Frontend}

Les fichiers \code{public/index.html} et \code{public/index_simple.html} ont été
mis à jour :
\begin{itemize}
  \item suppression de l'identité fixe \code{local-dev} pour les appels
  authentifiés;
  \item login sans \code{userRef};
  \item chat sans \code{userRef};
  \item panier sans \code{userRef} en query string;
  \item suggestions history sans \code{userRef};
  \item logout serveur via \code{/ai/logout}.
\end{itemize}

Un identifiant non sensible peut encore être utilisé pour du feedback
d'entraînement UI, mais il ne pilote plus les mutations panier ou commande.

\subsection{Réduction des logs sensibles}

Le service panier a été durci pour éviter de journaliser trop de détails :
\begin{itemize}
  \item pas de token;
  \item pas de mot de passe;
  \item pas d'adresse complète;
  \item pas de panier complet en debug générique;
  \item résumé limité : nombre de lignes, total, identifiant partiel, produit
  traité.
\end{itemize}

Cette réduction est importante pour la conformité : les logs sont souvent une
source secondaire de fuite de données personnelles.

\section{Pourquoi cette solution est plus sûre}

\subsection{Moindre confiance dans le client}

Le changement principal est architectural : l'identité utilisateur n'est plus un
paramètre métier envoyé par le navigateur. Elle devient une décision serveur
après authentification.

Cela empêche les scénarios suivants :
\begin{itemize}
  \item un client forge \code{userRef="local-dev"};
  \item un onglet réutilise une identité d'un autre test;
  \item deux comptes écrasent le même token runtime;
  \item une action panier s'applique au dernier compte connecté sous la même clé.
\end{itemize}

\subsection{Isolation par session}

Chaque login crée une session distincte. Deux navigateurs ou deux fenêtres
privées peuvent donc coexister :
\begin{itemize}
  \item le compte A a son cookie et son \code{web:<sessionId>};
  \item le compte B a un autre cookie et un autre \code{web:<sessionId>};
  \item les tokens Swissdata sont stockés sous deux clés différentes;
  \item panier et checkout suivent le token de la session courante.
\end{itemize}

\subsection{Extension WhatsApp passwordless}

Le canal WhatsApp ne peut pas s'appuyer sur le cookie \code{HttpOnly} utilisé
par l'interface web. Il a donc été traité comme un canal authentifié séparé,
avec son propre principal serveur \code{wa:<numero>}.

Le flux retenu est le suivant :
\begin{enumerate}
  \item le client écrit au numéro WhatsApp de l'assistant;
  \item Zelo AI normalise le numéro et construit le principal \code{wa:<numero>};
  \item si une session WhatsApp vérifiée existe déjà, la demande passe au moteur
  conversationnel;
  \item sinon, Zelo AI cherche via l'API Swissdata admin du shop le customer
  correspondant au numéro;
  \item si un customer unique est trouvé et lié à un user avec e-mail, Zelo AI
  génère un OTP à six chiffres;
  \item l'OTP est envoyé par e-mail et stocké côté Zelo AI uniquement sous forme
  de hash salé avec expiration;
  \item le client répond sur WhatsApp avec le code;
  \item après validation, Zelo AI crée une session serveur WhatsApp contenant
  le \code{customerId}, le \code{userId} et l'e-mail du compte vérifié;
  \item les commandes \code{add_items}, \code{show_cart}, \code{clear_cart},
  \code{discover_meals} et \code{exec_cart} utilisent alors exclusivement cette
  session.
\end{enumerate}

\subsection{Décision d'intégration avec Swissdata}

La revue du code Swissdata a montré que \code{/auth/token} supporte bien le
mode \code{double-auth}, mais seulement après création d'un token double-auth
par le processus d'authentification Swissdata existant. Swissdata n'expose pas,
via API publique, un endpoint permettant de créer un vrai token client à partir
du seul numéro de téléphone.

Pour respecter la contrainte de ne pas modifier \code{swissdata-api}, le flux
WhatsApp utilise donc un adapter \emph{customer-scoped} côté Zelo AI :
\begin{itemize}
  \item la preuve d'identité vient du code OTP envoyé à l'e-mail du user
  associé au numéro;
  \item les opérations panier passent par les routes Swissdata autorisées pour
  le shop;
  \item chaque action est filtrée sur le \code{customerId} vérifié;
  \item aucune session globale ou token partagé n'est réutilisé;
  \item la confirmation de commande WhatsApp est envoyée à l'e-mail du customer
  vérifié.
\end{itemize}

Cette approche évite l'accès direct Mongo en runtime, conserve le flux Meta
Cloud API existant, et sépare les identités WhatsApp sans demander de mot de
passe dans la conversation.

\subsection{Contrôles de sécurité WhatsApp}

\begin{longtable}{p{0.34\linewidth}p{0.58\linewidth}}
\toprule
\textbf{Contrôle} & \textbf{Mesure appliquée} \\
\midrule
Pas de mot de passe en chat & WhatsApp ne demande jamais le mot de passe Zelo. Les anciens messages \code{login ...} ne sont plus utilisés en mode OTP. \\
Anti-bruteforce & OTP à six chiffres, TTL court, nombre d'essais limité, cooldown avant renvoi. \\
Anti-replay & Le challenge est supprimé après validation, expiration ou verrouillage. \\
No account enumeration & Numéro inconnu, ambigu ou non lié : réponse générique, sans révéler si le compte existe. \\
Isolation session & Chaque session est indexée par \code{wa:<numero>} et liée au \code{customerId} vérifié. \\
Logs sûrs & Aucun OTP, token, mot de passe, hash, email complet, téléphone complet ou panier complet en logs. \\
Least privilege & En production, utiliser un compte Swissdata de service avec permissions minimales sur customers et orders du shop. \\
Routes legacy & Les routes \code{/link} sont admin-only en développement et désactivées en production. \\
\bottomrule
\end{longtable}

\section{Validation réalisée}

\subsection{Tests automatisés}

\begin{longtable}{p{0.34\linewidth}p{0.18\linewidth}p{0.40\linewidth}}
\toprule
\textbf{Test} & \textbf{Résultat} & \textbf{Objectif} \\
\midrule
\code{npm run test:session-isolation} & \statusok & Vérifier que deux sessions web gardent des tokens distincts et qu'une suppression ne contamine pas l'autre. \\
\code{npm run test:agent-formatting} & \statusok & Garantir que les réponses contrôlées restent valides. \\
\code{npm run test:whatsapp} & \statusok & Vérifier que le transport WhatsApp n'est pas cassé. \\
\code{npm run test:smoke} & \statusok & Vérifier le flux local complet : login, panier, add, clear, demandes longues. \\
\code{npm run build} & \statusok & Vérifier la compilation TypeScript. \\
\bottomrule
\end{longtable}

\subsection{Test multi-compte local}

Un scénario local a été exécuté avec deux comptes distincts dans
l'environnement Docker local :

\begin{enumerate}
  \item login du compte A;
  \item vidage du panier A;
  \item ajout d'un produit marqueur dans le panier A;
  \item login du compte B dans une session séparée;
  \item vidage du panier B;
  \item vérification que le panier B ne contient pas le produit du compte A;
  \item ajout de produits marqueurs dans le panier B;
  \item retour au compte A et vérification que les produits du compte B sont
  absents;
  \item envoi volontaire d'un payload contenant \code{userRef="local-dev"} avec
  le cookie du compte A;
  \item vérification que le backend ignore ce \code{userRef} forgé;
  \item passage de commande sur le compte B;
  \item nettoyage du panier A.
\end{enumerate}

Résultat : les deux paniers sont restés séparés, le checkout a utilisé la
session active correcte et la tentative de forger \code{local-dev} n'a pas
modifié le principal serveur.

\subsection{Ajout au preflight deploy}

Le script \code{deploy/01-local-preflight.sh} lance maintenant
\code{npm run test:session-isolation}. L'objectif est d'éviter une régression
avant un futur déploiement.

\section{Impact conformité nLPD/GDPR}

\subsection{Confidentialité}

Le correctif réduit fortement le risque qu'un utilisateur voie ou modifie le
panier d'un autre utilisateur. C'est directement lié au principe de
confidentialité et d'intégrité des données personnelles.

\subsection{Minimisation}

Le navigateur ne reçoit plus le token Swissdata dans la réponse de login. Il ne
porte qu'un cookie opaque, ce qui réduit la quantité de données sensibles
exposées côté client.

\subsection{Privacy by design}

La séparation des sessions est maintenant intégrée au design :
\begin{itemize}
  \item identité côté serveur;
  \item cookie \code{HttpOnly};
  \item logs redacted;
  \item suppression token/session au logout;
  \item tests dédiés anti-régression.
\end{itemize}

\subsection{Accountability}

La documentation, les tests et le preflight constituent des éléments de preuve
internes : ils montrent que le risque a été identifié, corrigé et validé.

\section{Limites et prochaines améliorations}

\subsection{Session store en mémoire}

Le store actuel est en mémoire Node.js. C'est acceptable pour un service
mono-instance, mais ce n'est pas suffisant pour un scaling horizontal ou un
redémarrage sans perte de session.

Amélioration recommandée \riskpone{} :
\begin{itemize}
  \item stocker les sessions dans Redis avec TTL;
  \item signer ou hacher les identifiants de session au repos;
  \item invalider proprement les sessions anciennes après rotation.
\end{itemize}

\subsection{Protection CSRF}

\code{SameSite=Lax} réduit le risque CSRF, mais une protection CSRF dédiée peut
être ajoutée si l'interface devient exposée publiquement avec cookies de
session.

Amélioration recommandée \riskpone{} :
\begin{itemize}
  \item token CSRF pour mutations web;
  \item validation d'origine \code{Origin}/\code{Referer} pour routes sensibles;
  \item rate limiting sur login/chat/checkout.
\end{itemize}

\subsection{Journalisation structurée}

Les logs sont plus sûrs, mais un vrai audit trail structuré serait préférable :
\begin{itemize}
  \item identifiant court de session;
  \item action métier;
  \item résultat;
  \item statut amont Swissdata;
  \item jamais de token, mot de passe, cookie ou panier complet.
\end{itemize}

\subsection{Tests production avant déploiement}

Avant le prochain déploiement, il faut refaire le scénario multi-compte sur
l'environnement cible :
\begin{itemize}
  \item deux navigateurs ou profils séparés;
  \item deux comptes Zelo distincts;
  \item paniers marqueurs différents;
  \item tentative de payload forgé;
  \item checkout sur un seul compte de test validé.
\end{itemize}

\section{Conclusion}

Le défaut initial était critique car il mélangeait l'identité client web avec
une valeur contrôlable et partagée. Le correctif appliqué remet l'identité au
bon endroit : côté serveur, après authentification, sous forme de session
opaque.

La solution améliore significativement la sécurité de Zelo AI :
\begin{itemize}
  \item séparation effective des sessions client;
  \item suppression de \code{local-dev} pour les mutations web;
  \item cookies \code{HttpOnly};
  \item routes sensibles protégées par session;
  \item logout serveur réel;
  \item logs moins sensibles;
  \item tests d'isolation ajoutés au preflight.
\end{itemize}

Le risque P0 identifié est donc traité pour l'architecture actuelle. Les
prochaines étapes concernent surtout le durcissement production : session store
Redis, protection CSRF/rate limit et audit trail structuré.

\end{document}
